\section{Introducción}\label{sec:intro}

El problema de las $n$ reinas pregunta cuántas maneras hay de colocar $n$ reinas en un tablero de
$n\times n$ sin que dos se ataquen. Denotamos ese número por $Q(n)$. En una solución, cada fila
contiene una reina; si anotamos su columna y leemos las filas de arriba abajo, obtenemos una
permutación, porque ninguna columna se repite. La Figura~\ref{fig:permutaciones} muestra por qué,
entre las $n!$ permutaciones posibles, sólo algunas son soluciones.

\begin{figure}[htbp]
\centering
\begin{tikzpicture}[x=5.6mm,y=5.6mm]
  \begin{scope}
    \bcuadricula{4}
    \foreach \r/\c in {1/2,2/4,3/1,4/3} {\breina{4}{\r}{\c}}
    \node[font=\small,anchor=north] at (2,-.35) {(a) Solución};
  \end{scope}
  \begin{scope}[shift={(9,0)}]
    \bcuadricula{4}
    \draw[btinta,densely dashed,line width=1.1pt] (.5,3.5)--(3.5,.5);
    \foreach \r/\c in {1/1,2/2,3/3,4/4} {\breina{4}{\r}{\c}}
    \node[font=\small,anchor=north] at (2,-.35) {(b) Diagonal compartida};
  \end{scope}
\end{tikzpicture}
\caption{Dos permutaciones en el mismo tablero. Las filas se numeran de arriba abajo y las columnas
de izquierda a derecha, empezando en uno. En (a) se leen las columnas $(2,4,1,3)$: las cuatro
reinas están también en diagonales distintas de cada orientación. En (b) se leen $(1,2,3,4)$:
cada fila y cada columna sigue teniendo una reina, pero la diagonal trazada contiene las cuatro.}
\label{fig:permutaciones}
\end{figure}


¿Cuánto reducen las diagonales esas $n!$ posibilidades cuando crece el tablero?
Simkin~\cite[Teorema~1.1]{Simkin} determinó el orden principal mediante un principio variacional:
existe una constante $\alpha$ tal que
\[
 Q(n)=\bigl((1+o(1))\,n\,e^{-\alpha}\bigr)^n .
\]
Nobel, Agrawal y Boyd~\cite{NAB} acotaron $\alpha$ en $[1{,}944000752,\,1{,}944001082]$.
Al tomar logaritmos y aplicar la fórmula de Stirling, el resultado de Simkin da
\[
 \log Q(n)-\log n!=n(1-\alpha)+o(n).
\]
El resto $o(n)$ no está determinado. La pregunta que motiva este artículo es la del término
siguiente: si, tras restar el término lineal, la diferencia converge, y hacia qué constante.

Construimos una cantidad determinista $Q_{\rm pred}(n)$, que llamamos predictor, cuyo término
constante sí podemos calcular. La construcción asigna un peso a cada casilla, de manera que cada fila
y cada columna sumen uno y ninguna diagonal reciba más de uno. Entre los pesos admisibles del modelo
se eligen los que maximizan la entropía de Bethe (Sección~\ref{sec:predictor}). Esa entropía aporta
el primer término del predictor. Los otros tres términos de la Definición~\ref{def:Qpred} son un
término local, que depende de los pesos por casilla; un pseudodeterminante, en el que interviene además cómo
se cortan las líneas; y una normalización. La definición exige que el maximizador sea interior, lo
que ocurre para todo $n$ suficientemente grande (Sección~\ref{sec:P02}).

Al pasar del tablero al cuadrado unidad aparece el perfil de masa uno $\nu_1$ de~\cite{Perfil}.
Escribimos $s_1$ para el valor óptimo en masa uno, que alcanza ese perfil, y $\Phi=s_1+1$. Bajo
la identificación con el problema de Simkin que se discute en la Sección~\ref{sec:preliminares}, y
que no usamos en ninguna prueba, $\Phi=1-\alpha$. La integral $I_2=\int\nu_1^2$ se toma sobre
el cuadrado unidad. Como $\nu_1$ tiene masa total uno, $I_2\ge1$, con igualdad sólo para la densidad
constante: $I_2$ mide cuánto se aparta el perfil de ella.

\begin{teorema}[Término constante del predictor]\label{thm:intro-constante}
Cuando $n\to\infty$,
\[
\begin{aligned}
 \log Q_{\rm pred}(n)-\log n!-n\Phi &\longrightarrow C_{\rm pred},\\
 C_{\rm pred} &=\tfrac12\log2-\tfrac12-\tfrac32I_2-\tfrac12\log{\det}_2(I+J),
\end{aligned}
\]
donde $J$ es el operador de Hilbert--Schmidt construido a partir de $\nu_1$ en la
Sección~\ref{sec:D03} e $I+J$ es definido positivo.
\end{teorema}

La prueba sigue la definición término a término. Primero se aproxima el maximizador discreto por
el perfil continuo con error $O(1/n)$ (Sección~\ref{sec:perfil}). Después se identifican por
separado las contribuciones de la entropía y del pseudodeterminante, y se suman al término local
(Sección~\ref{sec:constante}). El teorema se demuestra en la Sección~\ref{sec:D05}
(Teorema~\ref{thm:constante}).

El segundo resultado responde a una pregunta que ya se puede formular en un tablero pequeño.
Una colocación parcial deja un tablero residual $B$. Elegimos una fila vacía $r$ y llamamos
$B+c$ al residual que queda tras ocupar en ella una casilla legal $c$. Si $Q(B)$ cuenta las
completaciones de $B$, entonces
\[
 Q(B)=\sum_{c\in r} Q(B+c).
\]
En efecto, cada completación usa exactamente una de esas casillas. Una elección legal que luego
no pueda completarse aporta cero a la suma. La Figura~\ref{fig:residual} muestra esta partición
en un tablero de $5\times5$ con una reina en la fila~1 y la columna~2: en la segunda fila quedan dos
casillas legales, las de las columnas 4 y~5, y cada una da un sumando.

\begin{figure}[htbp]
\centering
\begin{tikzpicture}[x=5.6mm,y=5.6mm]
 \begin{scope}
  \bcuadricula{5}
  \foreach \r in {1,...,5} {\foreach \c in {1,...,5} {
    \pgfmathtruncatemacro{\ataque}{(\r==1)||(\c==2)||(abs(\r-1)==abs(\c-2))}
    \ifnum\ataque=1\relax\bmarca{5}{\r}{\c}\fi
  }}
  \breina{5}{1}{2}
  \foreach \c in {4,5}{\draw[btinta,line width=1.1pt] (\c-.5,3.5) circle[radius=.31];}
  \node[anchor=north,font=\small] at (2.5,-.32) {(a) Padre $B$};
 \end{scope}
 \draw[-{Stealth[length=2mm]},btinta] (5.55,2.5)--(7.35,2.5);
 \node[font=\footnotesize,above] at (6.45,2.65) {$(2,4)$};
 \begin{scope}[shift={(8.5,0)}]
  \bcuadricula{5}
  \foreach \r in {1,...,5}{\foreach \c in {1,...,5}{
    \pgfmathtruncatemacro{\ataque}{(\r==1)||(\c==2)||(abs(\r-1)==abs(\c-2))||(\r==2)||(\c==4)||%
(abs(\r-2)==abs(\c-4))}
    \ifnum\ataque=1\relax\bmarca{5}{\r}{\c}\fi
  }}
  \breina{5}{1}{2}\breina{5}{2}{4}
  \node[anchor=north,font=\small] at (2.5,-.32) {(b) Hijo $B+(2,4)$};
 \end{scope}
\end{tikzpicture}
\caption{La partición exacta del conteo en un tablero pequeño. En la segunda fila del tablero
izquierdo hay dos opciones legales, marcadas con aros. La flecha sigue una de ellas; la otra,
$(2,5)$, da el segundo sumando del conteo. Los puntos señalan casillas ocupadas o atacadas y las
coronas identifican las reinas colocadas. Las casillas disponibles conservan sus coordenadas.}
\label{fig:residual}
\end{figure}


El predictor asigna a $B$ el valor $e^{L(B)}$ (Sección~\ref{sec:B3T01}).
Dividimos la suma de las predicciones de los hijos por la predicción del padre:
\[
 \rho(B)=\frac{\sum_{c\in r}e^{L(B+c)}}{e^{L(B)}}
        =\sum_{c\in r}e^{L(B+c)-L(B)} .
\]
Si el predictor cumpliera exactamente la misma identidad que el conteo, este cociente valdría uno.
El siguiente teorema acota el defecto en escala logarítmica y especifica para qué residuales vale
esa cota.

Su enunciado usa nociones que se precisan en las Secciones~\ref{sec:insercion} y~\ref{sec:B5T01}.
Un residual con $k$ filas vacías tiene también $k$ columnas vacías y le faltan $k$ reinas; en
\cite{Completabilidad} ese número se denota $h$.
El \emph{punto de Bethe} de $B$ es el vector de pesos $p=(p_e)$, con un peso por casilla disponible
$e$, en el que se evalúa $L(B)$: un punto estacionario de la entropía de Bethe del residual bajo las
restricciones de fila y de columna. La \emph{masa} de una diagonal, de cualquiera de las dos
familias, es la suma de los pesos de sus casillas disponibles; su \emph{vacancia} es uno menos la
masa. El \emph{laplaciano de filas y columnas} es el del grafo cuyos vértices son las filas y las
columnas vacías y en el que cada casilla disponible $e$ es una arista de peso $p_e$ entre su fila y
su columna; su \emph{brecha} es su segundo autovalor más pequeño, contado con multiplicidad.

\begin{teorema}[Consistencia de segundo orden]\label{thm:intro-consistencia}
Sean $C,\eta,\gamma>0$ y sea $B$ un residual con $k$ filas vacías cuyo punto de Bethe $p$ es
estacionario e interior, con masas de fila y columna iguales a uno, $p_e\le C/k$, vacancias diagonales
al menos $\eta$ y brecha al menos $\gamma$ del laplaciano de filas y columnas ponderado por $p$.
\begin{enumerate}[(a)]
\item Para cada $\delta_0>0$ existen $k_0$ y $C_*$, que dependen sólo de $\delta_0,C,\eta,\gamma$, tales
que, si $k\ge\max(k_0,\delta_0n)$, entonces $|\log\rho(B)|\le C_*/k^2$ para toda fila vacía $r$
(Teorema~\ref{thm:macro} y Corolario~\ref{cor:rho}).
\item Existen $\varepsilon_0$, $k_0$ y $C_*$, que no dependen de $n/k$, tales que la misma conclusión
vale para $k\ge k_0$ sin la condición $k\ge\delta_0n$ si la masa de cada diagonal es a lo sumo
$\varepsilon_0$ (Teorema~\ref{thm:meso}).
\end{enumerate}
\end{teorema}

Ambas partes incluyen todos los hijos legales de la fila elegida. Para cada hijo se obtiene un
único punto estacionario en la rama pequeña, es decir, con todos sus pesos menores que $1/4$. La
primera parte pide que quede una fracción fija de las filas originales. La segunda prescinde de esa
condición cuando la masa de cada diagonal es pequeña. En los dos casos las hipótesis recaen sobre el
residual. El teorema es determinista: no afirma que los residuales que aparecen al colocar reinas al azar cumplan esas hipótesis.

Para trasladar el primer teorema al conteo exacto falta controlar una diferencia:
\[
 R_Q(n)=\log Q(n)-\log Q_{\rm pred}(n).
\]
Si $R_Q(n)\to0$, se obtiene
\[
 \log Q(n)=\log n!+n\Phi+C_{\rm pred}+o(1).
\]
No probamos que $R_Q(n)\to0$. La Sección~\ref{sec:discusion} explica qué estimaciones adicionales
sobre las trayectorias y los residuales con pocas filas vacías exigiría ese paso.

Numéricamente, $C_{\rm pred}\approx0{,}9401$: en todo el intervalo de $\alpha$ citado, los ajustes del predictor desde $n=40$ o $60$ dan entre $0{,}9400$ y $0{,}9402$, y una evaluación directa de la fórmula da $0{,}94014$ (Sección~\ref{sec:numerico}); ninguna de estas cifras está certificada.
Con la referencia $\Phi_0$ de esa sección, los conteos exactos para $16\le n\le27$ dan $\log Q(n)-\log n!-n\Phi_0$ entre $0{,}913$ y $0{,}941$.
